\subsection{实值函数族的包络}

定义 4. 设 \((f_{t})_{t\in I}\) 是定义在集合 X 上的一个实值函数族。在每一点 \(x\in X\) 处取值为 \(\sup_{t\in I}(f_t(x))\) {[}相应地，inf \((f_t(x))]\) 的 X 上的实值函数称为族 \((f_t)\) 的上包络（相应地，下包络），记作 \(\sup_{t\in I}f_{t}\) 或 \(\sup_{t}f_{t}\) （相应地，inf \(f_{t}\) 或 inf \(f_{t}\) ）。

族 \((f_{t})\) 的上包络就是这个族在定义于 X 上的实值函数的格 \(\overline{R}\) 中的上确界，这也就说明了记号 \(\sup f_{t}\) 的合理性。

此外，若给 \(\overline{\mathbf{R}}^{\mathbf{X}}\) 赋以其诸因子的拓扑（都与 \(\overline{R}\) 的拓扑相同）的乘积所成的拓扑，则我们有下述命题：

命题 10. 在积空间 \(\overline{\mathbf{R}}^{\mathbf{X}}\) 中，上包络 \(\sup f_{t}\) （实值函数族 \((f_{t})_{t\in I}\) 的上包络）是关于有向集 \(\mathfrak{F}(\mathbf{I})\) （由 \(\mathbf{I}\) 的有限子集组成）对映射 \(\mathrm{H}\to \sup_{t\in \mathbf{H}}f_t\) 取的极限{[}该映射把每个有限子集 \(\mathrm{H}\) （属于 \(\mathbf{I}\) ）映到有限子族 \((f_{t})_{t\in H}]\) 的上包络。

这由第 4 节的命题 5 以及第一章，§ 7，第 6 节，命题 10 的推论 1 立即得出。

因此我们可以写

\[
\sup _ {t \in I} f _ {t} = \lim _ {H \in \mathfrak {F} (I)} \left(\sup _ {t \in H} f _ {t}\right).\tag{9}
\]

定义 5. 实值函数族 \((f_t)_{i\in I}\) （定义在集合 \(\mathbf{X}\) 上）称为在 \(\mathbf{X}\) 中一致上有界（相应地，一致下有界）的，如果存在一个有限数 \(a\) ，使得 \(f_{t}(x) \leqslant a\) {[}相应地，\(f_{t}(x) \geqslant a\) {]} 对一切 \(x \in X\) 和一切 \(t \in I\) 成立。族 \((f_{t})\) 称为在 \(X\) 中一致有界的，如果它在 \(X\) 中既一致上有界又一致下有界。

于是，\((f_{t})\) 在 \(\mathbf{X}\) 中一致上有界当且仅当该族的上包络在 \(\mathbf{X}\) 中上有界。\((f_{t})\) 在 \(\mathbf{X}\) 中一致有界当且仅当族 \((|f_{t}|)\) 的上包络在 \(\mathbf{X}\) 中上有界{[}亦即当且仅当存在一个有限实数 \(a \geqslant 0\) ，使得 \(|f_{t}(x)| \leqslant a\) 对一切 \(x \in \mathbf{X}\) 和一切 \(t \in I\) 都成立 {]}。

